\section{Data Appendix}
\label{app:data}

\subsection{Variable Definitions}
\label{app:variables}

\begin{table}[htbp]
\centering
\caption{Variable Definitions}
\resizebox{\textwidth}{!}{
\begin{tabular}{lll}
\toprule
Variable & Construction & Source \\
\midrule
Book leverage & (dlttq + dlcq) / atq & Compustat fundq \\
Market leverage & (dlttq + dlcq) / (dlttq + dlcq + cshoq$\times$prccq) & Compustat fundq / CRSP \\
CF$_{it}$ & oibdpq / atq (winsorized 1\%--99\%) & Compustat fundq \\
Log assets & $\log(\texttt{atq})$ & Compustat fundq \\
Market-to-book & (cshoq $\times$ prccq) / ceqq & Compustat fundq \\
Profitability & oibdpq / atq & Compustat fundq \\
Tangibility & ppentq / atq & Compustat fundq \\
CF mean & 20-quarter rolling mean of CF$_{it}$ & Computed \\
CF volatility & 20-quarter rolling std.\ dev.\ of CF$_{it}$ & Computed \\
TRACE spread & (close\_yld $-$ GS10) $\times$ 100, in bps & WRDS TRACE / FRED \\
$\widehat{W}_{1,\text{shape}}$ & Eq.~\eqref{eq:w1_approx} and Section~\ref{sec:nu_identification} & Computed \\
\bottomrule
\end{tabular}
}
\end{table}

\subsection{CRSP-Compustat Merge}
\label{app:ccm}

I merge Compustat quarterly data to CRSP using the CRSP/Compustat Merged (CCM) link table (\texttt{crsp.ccmxpf\_lnkhist}). I retain observations where \texttt{linktype} $\in$ \{\texttt{LC}, \texttt{LU}, \texttt{LS}\} and \texttt{linkprim} $\in$ \{\texttt{P}, \texttt{C}\}, and require the fiscal quarter end to fall within the CCM link's effective date range (\texttt{linkdt} $\leq$ \texttt{datadate} $\leq$ \texttt{linkenddt}).

\subsection{Robustness of the Non-Parametric Representation}
\label{app:distribution_robustness}

The baseline specification represents $\hat{\mu}_{it}$ by ten empirical quantiles (decile midpoints). Results are robust to using five-point and twenty-point grids. \textbf{Cash flow measure}: the main measure is \texttt{oibdpq/atq} (quarterly EBITDA/assets), which is quarter-specific and requires no de-cumulation. Using de-cumulated \texttt{oancfy/atq} (quarterly operating cash flow from the cash flow statement) gives a within-firm coefficient of $-0.141^{***}$ ($t = -8.1$). Using \texttt{oibdpq/atq} with a log-normal parametric fit for $\hat{\mu}_{it}$ gives $-0.149^{***}$; using a gamma fit gives $-0.153^{***}$. \textbf{Distress controls} (Table~\ref{tab:robustness_distress}): adding the net income-to-assets ratio, a partial Whited-Wu financial constraints index \citep{WhitedWu2006financial}, and an EBIT-to-assets Z-score component \citep{Altman1968} jointly leaves the within-firm $\widehat{W}_{1,\text{shape}}$ coefficient at $-0.124^{***}$, confirming the within-firm pattern is not driven by financial distress in the conventional sense. \textbf{Capital supply robustness}: replacing the Compustat-based Fr\'{e}chet mean with a proxy anchored to FRED Flow-of-Funds aggregate financial assets (series \texttt{BOGZ1FL154090005Q}) yields $\rho = 0.87$ between the two $\hat{Q}_\nu$ functions and an unchanged within-firm leverage coefficient.

\begin{table}[htbp]
\centering
\caption{Robustness: $\widehat{W}_{1,\text{shape}}$ vs.\ Financial Distress Proxies}
\label{tab:robustness_distress}
\resizebox{\textwidth}{!}{
\begin{threeparttable}
\begin{tabular}{lcccc}
\toprule
 & (1) & (2) & (3) & (4) \\
\cmidrule(lr){2-5}
Outcome: Book leverage & Baseline & +NI/A & +WW & +All distress \\
\midrule
$\widehat{W}_{1,\text{shape}}$ & $-0.1619^{***}$ & $-0.1617^{***}$ & $-0.1245^{***}$ & $-0.1235^{***}$ \\
 & (0.0178) & (0.0178) & (0.0181) & (0.0181) \\[2pt]
CF mean & $-0.2118^{***}$ & $-0.2135^{***}$ & $-0.1984^{***}$ & $-0.1971^{***}$ \\
 & (0.0132) & (0.0132) & (0.0135) & (0.0136) \\[2pt]
CF volatility & $-0.3353^{***}$ & $-0.3331^{***}$ & $-0.3029^{***}$ & $-0.3029^{***}$ \\
 & (0.0194) & (0.0194) & (0.0195) & (0.0195) \\[2pt]
Log assets & $0.0185^{***}$ & $0.0191^{***}$ & $0.1828^{***}$ & $0.1828^{***}$ \\
 & (0.0008) & (0.0008) & (0.0086) & (0.0086) \\[2pt]
Market-to-book & $0.0034^{***}$ & $0.0033^{***}$ & $0.0035^{***}$ & $0.0035^{***}$ \\
 & (0.0001) & (0.0001) & (0.0001) & (0.0001) \\[2pt]
Profitability & $0.0187^{**}$ & $0.2085^{***}$ & $0.1127^{***}$ & $0.2356^{***}$ \\
 & (0.0065) & (0.0091) & (0.0100) & (0.0639) \\[2pt]
Tangibility & $0.1649^{***}$ & $0.1606^{***}$ & $0.1460^{***}$ & $0.1461^{***}$ \\
 & (0.0071) & (0.0071) & (0.0070) & (0.0069) \\[2pt]
Net income / assets &  & $-0.2006^{***}$ & $0.2497^{***}$ & $0.2499^{***}$ \\
 &  & (0.0062) & (0.0216) & (0.0216) \\[2pt]
WW index (partial) &  &  & $3.8883^{***}$ & $3.8874^{***}$ \\
 &  &  & (0.1978) & (0.1979) \\[2pt]
Z-score component (EBIT) &  &  &  & $-0.0396^{*}$ \\
 &  &  &  & (0.0205) \\[2pt]
\midrule
Firm FE & Yes & Yes & Yes & Yes \\
Year FE & Yes & Yes & Yes & Yes \\
$R^2$ & 0.0797 & 0.0848 & 0.1405 & 0.1405 \\
$N$ & 626,752 & 626,747 & 626,747 & 626,747 \\
\bottomrule
\end{tabular}
\begin{tablenotes}[flushleft]
\small
\item Dependent variable: book leverage (dlttq + dlcq) / atq. All specifications include firm and year-quarter fixed effects. Standard errors (in parentheses) clustered at firm level. NI/A = net income / total assets. WW index (partial) = $-0.091 \times ((\text{ibq}+\text{dpq})/\text{atq}) - 0.044 \times \log(\text{atq}) + 0.021 \times (\text{dlttq}/\text{atq})$, following \citet{WhitedWu2006financial}. Z-score component = $3.3 \times (\text{oibdpq}/\text{atq})$, the EBIT-to-assets term from \citet{Altman1968}. All variables winsorized at 1st/99th percentile. The baseline $N = 626{,}752$ (col.\ 1) exceeds the main regression sample $N = 625{,}105$ (Table~\ref{tab:panel_regressions}, col.\ 5) by 1,647 observations because the distress-robustness script applies a slightly less strict non-missing filter for the Frank-Goyal controls; within-firm coefficients on $\widehat{W}_{1,\text{shape}}$ are $-0.162$ in both samples (difference $< 0.001$). $^{***}$, $^{**}$, $^{*}$ denote significance at 1\%, 5\%, 10\%.
\end{tablenotes}
\end{threeparttable}
}
\end{table}

\textbf{Decile exclusion and tail trimming} (Table~\ref{tab:robustness_decile}): the within-firm coefficient on $\widehat{W}_{1,\text{shape}}$ becomes \emph{more negative} when decile 10 is excluded ($-0.382^{**}$ vs.\ $-0.163^{***}$ in the baseline) and under winsorization at p95 ($-0.499^{***}$) and p90 ($-0.607^{***}$). This pattern is theoretically informative, not a robustness failure. Decile-10 firms (mean $\widehat{W}_{1,\text{shape}} = 0.167$) have very low book leverage ($0.114$), consistent with the model's prediction, but their leverage is also constrained by non-distributional forces: distressed firms face equity issuance difficulties, covenant violations, and refinancing barriers that bound leverage from above \emph{and} below. Including them attenuates the linear regression coefficient because the leveraged-distress constraint limits the within-firm variation in leverage for the highest-$W_1$ firms. When these extreme-right-tail observations are removed, the relationship between distributional mismatch and leverage is estimated from a sample where the mapping between $\widehat{W}_{1,\text{shape}}$ and leverage is not confounded by distress-driven leverage constraints. The baseline ($-0.163$) is therefore a conservative lower bound on the distributional channel effect; the tail-trimmed estimates ($-0.382$ to $-0.607$) reflect the channel in the subpopulation of non-distressed firms.

\begin{table}[htbp]
\centering
\caption{Robustness: Sensitivity to $\widehat{W}_{1,\text{shape}}$ Outliers}
\label{tab:robustness_decile}
\resizebox{\textwidth}{!}{
\begin{threeparttable}
\begin{tabular}{lcccc}
\toprule
 & (1) & (2) & (3) & (4) \\
\cmidrule(lr){2-5}
Outcome: Book leverage & Baseline & Excl.\ decile 10 & Winsor.\ $\widehat{W}_{1,\text{shape}}$ (p95) & Winsor.\ $\widehat{W}_{1,\text{shape}}$ (p90) \\
\midrule
$\widehat{W}_{1,\text{shape}}$ & $-0.1625^{***}$ & $-0.3815^{**}$ & $-0.4992^{***}$ & $-0.6070^{***}$ \\
 & (0.0179) & (0.1351) & (0.0573) & (0.1125) \\[2pt]
\midrule
Firm FE & Yes & Yes & Yes & Yes \\
Year FE & Yes & Yes & Yes & Yes \\
$R^2$ & 0.0798 & 0.0883 & 0.0796 & 0.0788 \\
$N$ & 625,105 & 566,773 & 625,105 & 625,105 \\
\bottomrule
\end{tabular}
\begin{tablenotes}[flushleft]
\small
\item Dependent variable: book leverage. All specifications include firm and year-quarter fixed effects; standard errors clustered at firm level. Column (2) excludes firm-quarters in the top decile of $\widehat{W}_{1,\text{shape}}$. Columns (3)--(4) replace $\widehat{W}_{1,\text{shape}}$ with versions winsorized at the 95th and 90th percentile, respectively. $^{***}$, $^{**}$, $^{*}$ denote significance at 1\%, 5\%, 10\%.
\end{tablenotes}
\end{threeparttable}
}
\end{table}

\subsection{TRACE Bond Sample Construction}
\label{app:fisd}

TRACE trade-level data are accessed via WRDS (\texttt{trace.trade\_summary}), which provides the Dick-Nielsen (2009) cleaned daily corporate bond closing yields for the period 2012--2024. I apply the following filters: (1) bond type is corporate (excluding convertibles, asset-backed, and pass-throughs); (2) closing yield between 0.01\% and 30\%; (3) spread $= $ close\_yld $-$ GS10 between $-200$ and 2,850 bps; (4) at least five daily observations per firm-quarter. The CUSIP-to-Compustat match uses the six-digit issuer CUSIP from TRACE matched to \texttt{comp.names.cusip} from Compustat, retaining the first six characters of the nine-digit security CUSIP as the issuer CUSIP. After aggregation to firm-quarter means, the TRACE spread sample contains 56,837 firm-quarters for 2,019 unique firms.

\subsection{Trade-Off Theory as Special Case}
\label{app:tradeoff_special_case}

With a linear cost $c_{\mathrm{TO}}(x,d) = \kappa(d-x)_+$ (instead of quadratic) and a proportional tax benefit $\tau d$, the participation constraint reduces to: $\lambda \Prob_\mu(X < D^*) = \tau/(1+\tau)$, which is the standard trade-off optimality condition equating the marginal default probability to the marginal tax benefit. The quantile-coupling result $D^*(X) = F_\nu^{-1}(F_\mu(X))$ continues to hold because the linear shortfall cost $(d-x)_+$ is also submodular in $(x,d)$ and the comonotone coupling remains optimal. The Wasserstein distance under linear cost is $W_1(\mu,\nu)$, making the individual debt capacity $\delta_i^{\text{TO}} = W_1(\mu_i,\nu)/2$ formally identical to the quadratic case, though with different interpretations for the welfare formula ($W_1$ vs.\ $W_2^2$).

\subsection{Modigliani-Miller Invariance}
\label{app:mm_invariance}

Proved in the main appendix (Section~\ref{app:proofs}, subsection ``Microfoundation: Individual-Level Modigliani-Miller Invariance'').

\section{Estimation Details}
\label{app:estimation_details}

\subsection{Two-Stage Calibration Procedure}

The paper uses a two-stage calibration rather than full Simulated Method of Moments, for the following reason: the capital supply distribution $\hat{\nu}$ is identified directly from the cross-section of firm distributions as their Fr\'{e}chet mean (Section~\ref{sec:nu_identification}), without using any leverage or spread moments. This means the first stage has an analytical closed form and does not require simulation.

\paragraph{Stage 1: Calibrate $(\hat{\alpha}_\nu, \hat{\sigma}_\nu)$.}
Compute $\hat{Q}_\nu(u_k) = N^{-1}\sum_{it} \hat{q}_{it,k}$ for $k=1,\ldots,10$. Fit by OLS: $\log\hat{Q}_\nu(u_k) = \hat{\alpha}_\nu + \hat{\sigma}_\nu\Phi^{-1}(u_k) + \epsilon_k$. The estimate is analytical.

\paragraph{Stage 2: Calibrate $\hat{\kappa}$ via constrained GMM.}
Define the 10 moment conditions for $d = 1,\ldots,10$:
\begin{equation}
  g_d(\kappa) = \bar{s}_d^{\text{data}} - m_d(\kappa), \quad m_d(\kappa) \equiv \kappa\,x_d, \quad
  x_d = \frac{\bar{W}_{1,\text{shape},d}^{\text{spr}}}{2\hat{\mathbb{E}}_\nu[Y]} \times 10{,}000,
  \label{eq:gmm_moment}
\end{equation}
where $\bar{s}_d^{\text{data}}$ is the decile-mean TRACE spread and $\bar{W}_{1,\text{shape},d}^{\text{spr}}$ is the mean $\widehat{W}_{1,\text{shape}}$ of bond-issuing firms in decile $d$.

\textit{Optimal weighting matrix.} Estimate the $10\times 10$ covariance matrix $\hat{S}$ of the decile-mean spread vectors $(\bar{s}_1^{\text{data}},\ldots,\bar{s}_{10}^{\text{data}})$ from the 500 TRACE firm-level bootstrap replications: $\hat{S}_{dj} = \widehat{\mathrm{Cov}}(\bar{s}_d, \bar{s}_j)$ across replications. The optimal GMM weighting matrix is $\hat{W} = \hat{S}^{-1}$ (ridge-regularized for numerical stability).

\textit{GMM objective.} The constrained GMM estimator solves:
\begin{equation}
  \hat{\kappa} = \operatorname{argmin}_{\kappa \geq 0} \; Q(\kappa) = g(\kappa)'\,\hat{W}\,g(\kappa)
  \quad \text{subject to} \quad g_d(\kappa) \geq 0 \;\; \forall d.
  \label{eq:gmm_objective}
\end{equation}
The Jacobian (impulse function) of the model moments is $G = \partial m(\kappa)/\partial\kappa = x$ (analytic, since $m_d(\kappa) = \kappa x_d$ is linear in $\kappa$). The unconstrained GLS optimum is:
\begin{equation}
  \hat{\kappa}_{\text{WLS}} = \frac{x'\hat{W}\bar{s}^{\text{data}}}{x'\hat{W}x},
\end{equation}
with GLS asymptotic variance $\widehat{\mathrm{Avar}}(\hat{\kappa}_{\text{WLS}}) = (x'\hat{W}x)^{-1}$.

\textit{Constrained solution.} Empirically, $\hat{\kappa}_{\text{WLS}} = 0.142$ (Avar SE $= 0.004$) violates Proposition~\ref{prop:spreads}'s lower bounds for 2 deciles ($m_d(\hat{\kappa}_{\text{WLS}}) > \bar{s}_d^{\text{data}}$ at those deciles). Therefore the constrained optimum lies on the boundary of the feasible set: since $Q(\kappa)$ is quadratic and opens upward and the feasible set is $\kappa \leq \bar{s}_d/x_d$ for all $d$, the constrained minimizer is
\begin{equation}
  \hat{\kappa} = \min_{d\in\{1,\ldots,10\}}\;\frac{\bar{s}_d^{\text{data}}}{x_d}
               = \frac{\bar{s}_{d^*}^{\text{data}}}{x_{d^*}} = 0.104,
\end{equation}
where $d^* = 9$ is the binding decile. The constrained GMM objective at $\hat{\kappa}$ is $Q(\hat{\kappa}) = 594.8$; all 10 residuals $g_d(\hat{\kappa}) \geq 0$ by construction. The bootstrap SE for $\hat{\kappa}$ (500 TRACE firm-level replications of the full constrained GMM procedure) is $0.005$.

\subsection{Bootstrap Standard Errors}

Standard errors for $(\hat{\alpha}_\nu, \hat{\sigma}_\nu, \hat{\mathbb{E}}_\nu[Y])$ are computed via 500 Compustat firm-level bootstrap replications. Each replication resamples firms with replacement, re-computes the rolling-window empirical quantiles $\hat{q}_{it,k}$ on the bootstrap sample, re-computes the Fr\'{e}chet-mean $\hat{Q}_\nu(u_k)$, and re-fits $(\hat{\alpha}_\nu, \hat{\sigma}_\nu)$.

The bootstrap for $\hat{\kappa}$ (and for $\hat{S}$) uses 500 TRACE firm-level replications: each replication resamples TRACE bond-issuing firms with replacement, recomputes the 10 decile-mean TRACE spreads $(\bar{s}_1^{\text{data}},\ldots,\bar{s}_{10}^{\text{data}})$ and the mean $\widehat{W}_{1,\text{shape}}$ of bond issuers per decile, and re-solves the constrained GMM \eqref{eq:gmm_objective}. The sample covariance of the 500 decile-mean spread vectors gives $\hat{S}$ (used in the weighting matrix $\hat{W} = \hat{S}^{-1}$). The sample SD of the 500 $\hat{\kappa}$ values gives the bootstrap SE for $\hat{\kappa}$ ($= 0.005$). The two bootstrap procedures propagate, respectively, sampling uncertainty in the empirical cash flow quantiles (Stage~1) and in the TRACE spread moment conditions (Stage~2) through to the structural parameter estimates.
